\documentclass{article}
\usepackage[utf8]{inputenc}
\usepackage{polski}
\usepackage[polish]{babel}
\usepackage{amsmath}
\usepackage{amssymb}
\usepackage{geometry}

\geometry{a4paper, margin=1in}

\title{Teoria obliczeń i złożoność obliczeniowa -- Lista 11}
\author{Aleksandra Połacik}
\date{Październik 2025}

\begin{document}

\maketitle

\section{Zadanie 1}
Weryfikacja asymptotycznego tempa wzrostu funkcji przy użyciu definicji małego $o$. Sprawdzamy, czy dla $n \to \infty$ iloraz funkcji dąży do zera.

\begin{itemize}
    \item[(a)] $\sqrt{n} = o(n)$ \\
    Obliczamy granicę: $\lim_{n \to \infty} \frac{n^{1/2}}{n} = \lim_{n \to \infty} \frac{1}{n^{1/2}} = 0$. Warunek spełniony.

    \item[(b)] $n = o(n \cdot \log \log n)$ \\
    Skracamy przez $n$: $\lim_{n \to \infty} \frac{n}{n \log \log n} = \lim_{n \to \infty} \frac{1}{\log \log n} = 0$, ponieważ mianownik rośnie nieograniczenie.

    \item[(c)] $n \log \log n = o(n \log n)$ \\
    Po uproszczeniu otrzymujemy $\lim_{n \to \infty} \frac{\log \log n}{\log n}$. Stosując podstawienie $t = \log n$, mamy $\lim_{t \to \infty} \frac{\log t}{t} = 0$.

    \item[(d)] $n \log n = o(n^2)$ \\
    Badamy iloraz: $\lim_{n \to \infty} \frac{n \log n}{n^2} = \lim_{n \to \infty} \frac{\log n}{n} = 0$. Logarytm rośnie wolniej niż funkcja liniowa.

    \item[(e)] $n^2 = o(n^3)$ \\
    Granica $\lim_{n \to \infty} \frac{n^2}{n^3} = \lim_{n \to \infty} \frac{1}{n} = 0$. Wielomian niższego stopnia jest małym $o$ wielomianu wyższego stopnia.
\end{itemize}

\section{Zadanie 2}
Należy wykazać, że dla funkcji dodatniej $f(n)$ nie może zachodzić $f(n) = o(f(n))$.

Z definicji małego $o$, relacja ta zachodziłaby, gdyby $\lim_{n \to \infty} \frac{f(n)}{f(n)} = 0$. Jednak dla dowolnej funkcji $f: \mathbb{N} \to \mathbb{R}^+$, iloraz $\frac{f(n)}{f(n)}$ jest zawsze równy $1$. Granica stałej wartości $1$ przy $n$ dążącym do nieskończoności wynosi $1$, a nie $0$. Ponieważ $1 \neq 0$, stwierdzenie $f(n) = o(f(n))$ jest zawsze fałszywe.

\section{Zadanie 3}
Analiza prawdziwości stwierdzeń dotyczących notacji dużego $O$:

\begin{itemize}
    \item[(a)] $2n = O(n)$ -- \textbf{PRAWDA}. Stałe współczynniki nie wpływają na klasę złożoności (można przyjąć $c=2$).
    \item[(b)] $n^2 = O(n)$ -- \textbf{FAŁSZ}. $n^2$ rośnie istotnie szybciej niż funkcja liniowa.
    \item[(c)] $n^2 = O(n \log^2 n)$ -- \textbf{FAŁSZ}. Kwadrat $n$ dominuje nad iloczynem $n$ i logarytmów
    \item[(d)] $n \log n = O(n^2)$ -- \textbf{PRAWDA}. Funkcja $n \log n$ jest ograniczona z góry przez $n^2$ dla $n \ge 1$.
    \item[(e)] $3^n = 2^{O(n)}$ -- \textbf{PRAWDA}. Ponieważ $3^n = 2^{n \log_2 3}$, a $n \log_2 3$ mieści się w $O(n)$.
    \item[(f)] $2^{2^n} = O(2^{2^n})$ -- \textbf{PRAWDA}. Relacja zwrotna zawsze zachodzi dla dużego $O$.
\end{itemize}

\section{Zadanie 4}
Weryfikacja relacji małego $o$:

\begin{itemize}
    \item[(a)] $n = o(2n)$ -- \textbf{FAŁSZ}. Granica ilorazu wynosi $1/2 \neq 0$.
    \item[(b)] $2n = o(n^2)$ -- \textbf{PRAWDA}. Iloraz $2/n$ zbiega do zera.
    \item[(c)] $2n = o(3^n)$ -- \textbf{PRAWDA}. Każda funkcja wykładnicza dominuje nad wielomianową.
    \item[(d)] $1 = o(n)$ -- \textbf{PRAWDA}. Funkcja stała rośnie wolniej niż liniowa.
    \item[(e)] $n = o(\log n)$ -- \textbf{FAŁSZ}. Jest dokładnie odwrotnie; $n$ rośnie szybciej.
    \item[(f)] $1 = o(1/n)$ -- \textbf{FAŁSZ}. Granica $\frac{1}{1/n} = n \to \infty$.
\end{itemize}

\section{Zadanie 5}
Wyznaczenie maksymalnego rozmiaru danych $n$, które można przetworzyć w zadanym czasie $t$, przy założeniu, że operacja trwa $f(n)$ mikrosekund.

Wartości czasu w mikrosekundach ($\mu s$):
\begin{itemize}
    \item 1 s = $10^6 \mu s$
    \item 1 min = $6 \cdot 10^7 \mu s$
    \item 1 h = $3.6 \cdot 10^9 \mu s$
    \item 1 dzień = $8.64 \cdot 10^{10} \mu s$
    \item 1 rok = $3.15 \cdot 10^{13} \mu s$
\end{itemize}

\begin{table}[h]
\centering
\begin{tabular}{|l|c|c|c|c|c|}
\hline
$f(n)$ & 1 sekunda & 1 minuta & 1 godzina & 1 dzień & 1 rok \\ \hline
$\log_2 n$ & $2^{10^6}$ & $2^{6 \cdot 10^7}$ & $2^{3.6 \cdot 10^9}$ & $2^{8.6 \cdot 10^{10}}$ & $2^{3.1 \cdot 10^{13}}$ \\ \hline
$\sqrt{n}$ & $10^{12}$ & $3.6 \cdot 10^{15}$ & $1.3 \cdot 10^{19}$ & $7.4 \cdot 10^{21}$ & $9.9 \cdot 10^{26}$ \\ \hline
$n$ & $10^6$ & $6 \cdot 10^7$ & $3.6 \cdot 10^9$ & $8.6 \cdot 10^{10}$ & $3.1 \cdot 10^{13}$ \\ \hline
$n^2$ & $10^3$ & $7745$ & $60\,000$ & $293\,938$ & $5\,615\,692$ \\ \hline
$2^n$ & $19$ & $25$ & $31$ & $36$ & $44$ \\ \hline
$n!$ & $9$ & $11$ & $12$ & $13$ & $16$ \\ \hline
\end{tabular}
\end{table}

\section{Zadanie 6}
Dowód przechodniości w notacji $O$.
Założenia:
\begin{enumerate}
    \item $f(n) \le c_1 \cdot g(n) \cdot h_1(n)$ dla $n \ge n_1$
    \item $h_1(n) \le c_2 \cdot h(n)$ dla $n \ge n_2$
\end{enumerate}
Podstawiając (2) do (1) dla $n \ge \max(n_1, n_2)$:
\[ f(n) \le c_1 \cdot g(n) \cdot (c_2 \cdot h(n)) = (c_1 \cdot c_2) \cdot (g(n) \cdot h(n)) \]
Przyjmując nową stałą $C = c_1 \cdot c_2$, otrzymujemy $f(n) \in O(g(n) \cdot h(n))$, co kończy dowód.

\end{document}